\chapter{Modern Direct3D}
\label{ch:modern-direct3d}
\label{ch:direct3d-backends}

Four CNA identities occupy the modern Direct3D line.  Three broad implementations share one
development loop; \texttt{DIRECTX10} is a deliberately narrower bridge between them:
cross-compiled from this project's own Debian machine using the same MinGW-w64 toolchain
Chapter~\ref{ch:building-cna}'s Windows cross-build section already introduced, then run and
verified under Wine with GPU execution --- \texttt{D3D9} and \texttt{D3D11} through DXVK,
\texttt{D3D12} through a different translation layer, vkd3d-proton, in its own separate Wine
prefix. Native-Windows CI now covers \texttt{DIRECTX11} and \texttt{DIRECTX12}; local Wine
evidence remains a distinct tier and \texttt{DIRECTX9}/\texttt{DIRECTX10} have no equivalent CI
job.

\begin{longtable}{p{0.15\textwidth}p{0.21\textwidth}p{0.22\textwidth}p{0.24\textwidth}}
\toprule
Identity & Shader model & Extended route & RT / MRT / query \\
\midrule
\endfirsthead
\toprule
Identity & Shader model & Extended route & RT / MRT / query \\
\midrule
\endhead
\texttt{DIRECTX9} & D3D9 / XNA-native effects & native & yes / yes / yes \\
\texttt{DIRECTX10} & two embedded SM4 shaders & inherited colored downgrade & yes / yes / no \\
\texttt{DIRECTX11} & SM5 shared generated package & native & yes / yes / yes \\
\texttt{DIRECTX12} & SM5 plus PSO/root-signature caches & native & yes / yes / yes \\
\bottomrule
\end{longtable}

\section{D3D9: zero-tolerance XNA oracle comparisons}

\texttt{D3D9} targets pixel-for-pixel agreement with XNA~4.0 on its measured corpus. It uses
Microsoft's vendored XNA stock-effect HLSL and compiles it through
\texttt{d3dcompiler\_47.dll}; it does not substitute a reimplemented shader package.

\subsection{Oracle: XNA~4.0 under Wine}

A dedicated tool, \texttt{tools/xna-oracle/}, runs XNA~4.0 under Wine using the GAC assemblies
and the in-prefix \texttt{csc.exe}. It renders each declarative scene through XNA and CNA's
\texttt{D3D9} renderer, then compares the images at zero tolerance. Both sides use the same DXVK layer on
the same GPU, isolating CNA behavior from differences between graphics stacks. The pinned corpus contains
\XnaOracleSceneCount{} scene files. Every retained comparison is an exact match under this
Wine/DXVK configuration; the result does not cover paths absent from that corpus.

\subsection{Current contract and evidence}

Full device lifecycle via plain \texttt{Direct3DCreate9} (not D3D9Ex, a deliberate choice),
with \texttt{DeviceLost} / \texttt{DeviceResetting} / \texttt{DeviceReset} handling; all
five XNA stock effects, with 61 of 66 compiled shader variants byte-identical to Microsoft's
own shipped bytecode (the remaining five differ only by HLSL compiler version, and are
separately compared through the oracle); an
oracle- and mutation-verified \texttt{SpriteEffect.fx}, including the half-pixel offset
Direct3D~9 is notorious for --- caught before it could become a false-positive ``closed''
claim, because the offset's absence is invisible on a $1\times1$ test texture and only shows
up at a larger size; and \cnaclass{GraphicsProfile.Reach} / \texttt{.HiDef} enforcement
through a \texttt{D3DCAPS9} query --- the only audited CNA renderer
where the distinction reaches a native capability query,
as Chapter~\ref{ch:graphicsdevice} already noted. Render-target sampling and the
\texttt{PreferPerPixelLighting}/specular gaps described in Chapter~\ref{ch:stock-effects}
were corrected. The specular repair produced zero-tolerance frame matches for four of five
\texttt{PixelLighting} bytecode variants via the
oracle. The fifth, an untextured vertex-lit variant, is unreachable at the pin because the
required position-only vertex layout is missing on the measured stock-effect paths.

\begin{sourcenote}
D3D9's lost-device lifecycle is reachable through two framework-reserved keys:
after ordinary input processing, non-repeated F9 enters \texttt{Lost} and raises
\texttt{DeviceLost}; F10 performs the reset, bracketing it with
\texttt{DeviceResetting} and \texttt{DeviceReset}. The smoke suite checks exact event counts,
that rendering fails while lost, exact pixels after recovery, and release/recreation of
default-pool resources. This strong deterministic debug path should not be confused with two
other public extensions: D3D9's \texttt{SetStringMarkerEXT()} throws, and
\texttt{SetContextRecoveryEnabled()} changes the shared \cnaclass{Texture2D} shadow policy
before its D3D9 hook throws. Section~\ref{sec:backend-recovery-debug-matrix} compares all
fourteen products in the inherited audit cohort and calls out that partial-state outcome.
\end{sourcenote}

\subsection{Coverage limits}

Every \cnaclass{SurfaceFormat} besides \texttt{Color} remains unsupported (a shared,
cross-renderer limitation from Chapter~\ref{ch:textures-rendertargets}, not a D3D9-specific
one). \texttt{Texture3D} has no oracle-scene coverage, but it has a capability-gated
sub-volume \texttt{SetData()}/\texttt{GetData()} byte round trip: \texttt{D3D9\_Smoke} creates a
\(4\times4\times4\) volume only when \texttt{D3DCAPS9::MaxVolumeExtent}
is nonzero, writes 32 distinct RGBA bytes into the off-origin \(2\times2\times2\) box at
\((1,1,1)\), and requires an exact readback of that same box. The test proves D3D9's
\texttt{LockBox()} row- and slice-pitch copy in both directions; it does not prove an oracle
rendering path or automatic mip generation. \cnaclass{OcclusionQuery} is not
built for this renderer at all, the same gap D3D12 carries; and
\cnaclass{SpriteSortMode.Immediate} / \texttt{.Texture} were confirmed not viable as
deterministic oracle scenes at all --- \texttt{Immediate} is not something a raster diff can
observe, and \texttt{Texture} sorts by an implementation-defined \texttt{GetHashCode()},
which no two implementations are obligated to agree on. The caveat that applies to every
result in this section: \texttt{D3DCAPS9} itself is synthesized by DXVK on this dev loop, not
reported by XNA-era Windows driver hardware --- native-Windows hardware verification remains
open, as stated above.

\subsection{D3D9 render-target limits and evidence}

The D3D9 smoke test establishes render-target storage and recovery. An 8x8 target with
Depth24Stencil8 is bound, cleared
through the public device, and copied with D3D9's required
\texttt{GetRenderTargetData()}--system-memory-surface--\texttt{LockRect()} route; its own surface
must contain the exact requested color. A cube-face target has the same proof for one selected
face, followed by an independent back-buffer clear/readback after unbind.

MSAA is capability-gated by \texttt{IDirect3D9::CheckDeviceMultiSampleType()}, not assumed from
the requested count. When 4x is available, a multisampled target is cleared and unbound; D3D9's
\texttt{StretchRect()} resolve must make its separate sampleable texture read exactly
\((200,210,220,255)\). The assertion reads the resolved result, not the multisample surface.
One naming trap matters here: this is an internal device-capability check in the creation path,
not \cnaclass{GraphicsDevice}\texttt{::SupportsCapability}. D3D9 never overrides that public
method, so \texttt{MultiSampleAntiAliasing} still reports true even when
\texttt{CheckDeviceMultiSampleType()} rejects every count above one. The same mismatch affects
\texttt{AnisotropicFiltering}: sampler creation uses \texttt{D3DTEXF\_ANISOTROPIC} and
\texttt{D3DSAMP\_MAXANISOTROPY}, but neither the public query nor the setter consults and clamps
to \texttt{D3DCAPS9::MaxAnisotropy}.

Multiple render targets are bounded by \texttt{D3DCAPS9::NumSimultaneousRTs}: two
targets receive one exact-color clear, unbinding restores the back buffer, and requesting one
more target than the device reports throws rather than silently truncating. This proves D3D9's
multi-attachment binding and clear semantics. It does not separately prove a shader with several
fragment outputs, so the result should not be generalized into a G-buffer draw claim.

\subsection{One oracle scene}

The oracle's own scene format is a small, declarative text file, not a compiled asset or a
piece of C\texttt{++}. The simplest scene in the corpus, \texttt{colored3d.scene}, fits in
eleven lines:

\begin{lstlisting}
width=256
height=256
profile=HiDef
clearcolor=100,149,237,255
vertexcolor=true
lighting=false
primitive=TriangleList
vertex=-0.6,-0.6,0,255,0,0,255
vertex=0.0,0.7,0,0,255,0,255
vertex=0.6,-0.6,0,0,0,255,255
\end{lstlisting}

With \texttt{vertexcolor=true} and \texttt{lighting=false}, \cnaclass{BasicEffect} draws one
red/green/blue-cornered triangle under identity transforms over
\texttt{Color.CornflowerBlue}, RGBA $(100,149,237,255)$. XNA and CNA render the scene
independently. Matching corners establish vertex-color plumbing; matching interior pixels also
exercise Gouraud interpolation. The more elaborate \texttt{fog\_gradient\_quad}
(Chapter~\ref{ch:easygl-backend}'s own negative-\texttt{FogEnd} finding), uses the identical
file format with a handful of additional \texttt{fog*} keys --- the same declarative scheme
scales from this eleven-line triangle through the corpus's \XnaOracleSceneCount{} scenes.

\subsection{Two defects exposed by the oracle}

The first was a sort-mode Z-clipping defect: \texttt{zFarPlane=1}
mapped a nonzero \cnaclass{SpriteBatch} layer depth outside Direct3D~9's own valid clip-space
Z range, silently clipping away any sprite with a nonzero layer depth entirely --- fixed by
using \texttt{zFarPlane=-1} instead. The second was shared across renderers:
\cnaclass{GraphicsProfile} never reached the device because \texttt{Game}'s
\cnaclass{GraphicsDevice} was eagerly default-constructed
(hardcoded to \texttt{Reach}) before \cnaclass{GraphicsDeviceManager} even existed to apply a
game's requested profile --- meaning a game setting \texttt{GraphicsProfile = HiDef;
ApplyChanges();} had no path to the live device at all. This is a shared, cross-renderer bug,
only \emph{discovered} while building \texttt{D3D9}'s own tests, because \texttt{D3D9} is the
one renderer where the profile actually matters enough to notice.

\subsection{Why the \texttt{zFarPlane} sign hid nonzero-depth sprites}

The sign controls clip-space placement. \cnaclass{SpriteBatch}'s orthographic projection maps a
\texttt{layerDepth} value in $[0, 1]$ (XNA's own documented range: 0 = front, 1 = back) onto
Direct3D~9's own native clip-space Z range, which is $[0, 1]$ --- \emph{not} the
$[-1, 1]$ range OpenGL-family APIs use. Constructing that projection with
\texttt{zFarPlane=1} produces a near/far mapping where any \texttt{layerDepth} strictly
greater than what the incorrect matrix treats as its own valid far bound gets clipped by the
GPU's fixed-function Z-clip test without an error. The default \texttt{layerDepth=0} was
unaffected, so the defect appeared only with nonzero depth ordering. The
\texttt{zFarPlane=-1} matrix maps the
full documented $[0,1]$ \texttt{layerDepth} range maps inside Direct3D~9's own valid clip
volume instead of partly outside it.

\section{Three D3D11 resource-lifetime groups}
\label{sec:d3d11-lifetime-groups}

Window resize and device-removed recovery affect different resource sets. D3D11 separates them
into three lifetime groups:

\begin{description}
  \item[Device lifetime.] \texttt{ID3D11Device}, \texttt{ID3D11DeviceContext}, the
        \texttt{IDXGIFactory2} chain, and a cached \texttt{allowTearingSupported\_} capability
        flag --- created once at startup, torn down only by full device-removed recovery.
  \item[Swap-chain lifetime.] The \texttt{IDXGISwapChain1} object itself --- also created once
        at startup (and recreated only alongside the device, on recovery), but a plain resize
        deliberately \emph{reuses the same object} via
        \texttt{IDXGISwapChain1::ResizeBuffers(\ldots{})} rather than destroying and
        recalling \texttt{CreateSwapChainForHwnd}.
  \item[Window-size lifetime.] The back-buffer \texttt{ID3D11RenderTargetView}, the
        depth-stencil texture and its \texttt{ID3D11DepthStencilView}, and the viewport ---
        unbound and released \emph{before} \texttt{ResizeBuffers}, then recreated
        \emph{after} it, on every resize and on device-removed recovery alike.
\end{description}

Resize releases and recreates only the window-size group, with one \texttt{ResizeBuffers()} call
on the existing swapchain. Device removal rebuilds all three groups. This division avoids both
recreating the device on resize and omitting the swapchain from recovery.

All D3D11, DXGI, and D3D12 interfaces use \texttt{Microsoft::WRL::ComPtr<T>}. A MinGW-w64/Wine
spike exercised device creation, \texttt{.As()}, \texttt{GetParent()}, and the
\texttt{ReleaseAndGetAddressOf()} repopulation needed by resize before that ownership model was
adopted.

\section{DIRECTX10: a bounded MVP with one unsafe default}

\texttt{DIRECTX10} uses Wine's D3D10/D3D10.1 forwarding path into DXVK's
\texttt{d3d10core} and DXGI.  Its two embedded shaders provide colored 2D and 3D submission,
render targets, and MRT, but it deliberately leaves effect creation, occlusion queries,
Texture3D, TextureCube, and cube targets at their base null implementations.

The sharpest boundary is not a throw. The renderer does not override either effect-aware
extended primitive method, so the base class calls its colored draw implementation after
discarding \cnaclass{GpuDrawParams}.  A BasicEffect-shaped call can therefore produce pixels
while silently losing texture, lighting, fog, skinning, PBR, instancing, or custom-effect
intent. This is the sole unconditional inherited colored downgrade among all
\RendererFamilyCount{} renderer
families.

That behavior makes capability checks necessary but insufficient. \texttt{ThreeD} can be true
because colored geometry is implemented while effect semantics remain absent. A porter should
treat \texttt{DIRECTX10} as an explicit MVP target and verify the exact draw vocabulary used by
the game rather than infer parity from a successful triangle.

\section{D3D11: native Windows implementation}

\texttt{D3D11} is CNA's first renderer built directly on Direct3D~11, verified in the
same MinGW-w64-plus-Wine-plus-DXVK loop as \texttt{D3D9}. Its covered path includes an
\texttt{ID3D11Device} / \texttt{ID3D11DeviceContext} pipeline: buffers, textures, and render
targets (including MSAA and multiple render targets), state objects, all ten stock HLSL
shader variants across the five stock effects, \cnaclass{SpriteBatch}, and --- unique
to this renderer and \texttt{D3D12} --- a runtime-\texttt{D3DCompile()}-backed
\texttt{ShaderEffect} path for custom shaders, already introduced in
Chapter~\ref{ch:shader-fx-gap}. Pixel assertions use GPU readback, so they establish more than
successful API return codes.

Device-removed recovery has \emph{detection-only evidence}: the check is wired, but the
Wine/DXVK campaign did not trigger device removal, so full recovery remains unobserved.

An early build exposed a static-link cycle when the linker could not resolve
\texttt{Effect::Apply()}: the D3D11 \cnaclass{SpriteBatch}
renderer calls back into \texttt{Effect::Apply()}, which lives in the main CNA library, which
itself depends on the renderer). MinGW's single-pass archive resolution never revisited the
already-searched \texttt{libCNA.a} to satisfy it. The fix was an explicit repeated
\texttt{target\_link\_libraries(\ldots{} CNA)} for the D3D11/D3D12 targets, using CMake's
documented support for a repeated static-library dependency. The failure reproduced on the base
revision and was not introduced by the D3D11 work.

\subsection{The public presentation-mode call reaches D3D11 and D3D12, but neither scales yet}

Chapter~\ref{ch:game-loop}'s \texttt{CNAEXT} \texttt{Presentation\allowbreak{}Mode} is not lost before it
reaches these renderers: the public manager/device forwarding path passes the selected enum to each
renderer. But the final handlers are
currently intentional no-ops. D3D11 explicitly discards \texttt{mode}; D3D12 has the same
empty handler. They create and present their swap chain at the SDL window's pixel size,
not a separately rendered logical surface which could be letterboxed, cropped, or stretched.

Their \texttt{SetVirtualResolution(width,height)} methods only retain the supplied integers.
On D3D11, \texttt{GetViewportSize()} still queries the window's physical pixel size; on D3D12,
the stored virtual dimensions are returned to the caller but do not create a scaling pass.
Consequently setting \texttt{Letterbox}, \texttt{Overscan}, \texttt{Stretch},
\texttt{NativeBackBuffer}, or \texttt{FixedHeightDynamicWidth} produces no different
presentation geometry on either native Direct3D renderer today. This is a renderer limitation,
not an XNA-compatibility promise: the enum itself is a CNA extension, and there is no test that
could honestly pixel-prove a mode-specific result until a logical render target plus final
composite/scaling path exists.

\subsection{One \texttt{PresentInterval::Two}, two opposite Direct3D mistakes}

The three renderers do not share one presentation-interval policy. D3D9 preserves CNA's
0/1/2 conversion in \texttt{D3DPRESENT\_}\allowbreak\texttt{PARAMETERS}: Immediate becomes
\texttt{D3DPRESENT\_}\allowbreak\texttt{INTERVAL\_}\allowbreak\texttt{IMMEDIATE}, Two becomes
\texttt{D3DPRESENT\_}\allowbreak\texttt{INTERVAL\_}\allowbreak\texttt{TWO}, and everything
else becomes \texttt{D3DPRESENT\_}\allowbreak\texttt{INTERVAL\_}\allowbreak\texttt{ONE}.
A runtime change marks presentation state dirty and
immediately calls the same device-reset path used for a resize. This is the only CNA renderer
in the Direct3D family whose source requests a literal native two-retrace interval.

It is not yet a safe implementation. Microsoft's D3D9 contract says interval Two must be
present in \texttt{D3DCAPS9::PresentationIntervals}, and windowed mode supports only Default,
Immediate, and One.\footnote{\url{https://learn.microsoft.com/en-us/windows/win32/direct3d9/d3dpresent}}
CNA checks neither condition. Its ordinary game window is windowed, so requesting Two can
make \texttt{IDirect3DDevice9::Reset()} return failure. The renderer logs that result, restores
its previous dimensions, leaves \texttt{presentationDirty\_} true, and retries at each later
\texttt{Present()}; the public \cnaclass{PresentationParameters} still reports Two. No
dedicated test exercises this route, and DXVK's synthesized capability structure would not
replace the still-needed real-Windows check anyway.

D3D11 and D3D12 make the opposite mistake. Both reduce the construction argument and every
later setter call to \texttt{vsyncEnabled\_ = interval > 0}; their
\texttt{Present()} then passes only sync interval~1 or~0 to DXGI. DXGI itself accepts
1 through~4 and defines them as waiting for at least the corresponding number of vertical
blanks,\footnote{\url{https://learn.microsoft.com/en-us/windows/win32/api/dxgi/nf-dxgi-idxgiswapchain-present}}
so nothing in the native API forces CNA to collapse Two. The real windowed smoke paths prove
that default presentation succeeds, but the D3D plans explicitly say the no-VSync/tearing
branch is untested; no located test distinguishes sync interval~2 either.

The useful porting conclusion is deliberately narrow. D3D9 has the correct representation
but incomplete validation and failure reporting; D3D11/12 have robust per-present selection
but an unnecessarily Boolean representation. A successful reset or a stored
\texttt{PresentInterval::Two} is not yet evidence of half-refresh pacing on any of the three.
The cross-renderer table in \S\ref{sec:backend-swap-interval-matrix} places these results beside
the other eleven products in that audit cohort.

\subsection{Only D3D9 consumes the presentation format/fullscreen hook}

D3D9 likewise stands apart for the other three presentation fields. It maps the requested
backbuffer enum to \texttt{D3DFORMAT} (with a required Color-specific A8B8G8R8-to-A8R8G8B8
display substitution), maps None/Depth16/Depth24/Depth24Stencil8 to no attachment,
D16, D24X8, and D24S8 respectively, and writes
\texttt{D3DPRESENT\_PARAMETERS.Windowed = !IsFullScreen}. A later manager reset reaches
\texttt{UpdatePresentationFormatEXT()} and immediately resets the native device. The smoke
suite observes a real D24S8 surface after manager-driven setup, so this is more than a stored
field.

D3D11 and D3D12 ignore all three renderer arguments. Their windowed swapchains are fixed
R8G8B8A8 UNORM and their default depth attachments fixed D24S8, even when public state says
Depth24 or None. Shared SDL code may still enlarge the containing window for fullscreen.
Both \texttt{CreateSwapChainForHwnd} calls omit a fullscreen descriptor, and neither renderer
calls DXGI \texttt{SetFullscreenState}; their \texttt{exclusiveFullscreen\_} fields remain false.
D3D12's \texttt{HeadlessEXT} route is a third case: no window means no swapchain and no
default depth attachment at all. It can render only into explicitly-created targets.

These distinctions prevent two tempting overclaims. D3D9's enum mapping is not a guarantee
that every mapped format is a legal display mode; failed native reset can leave the previous
attachments active while public state has already changed. And an SDL fullscreen window on
D3D11/12 is not DXGI exclusive fullscreen. The inherited fourteen-product comparison and test boundary
are in \S\ref{sec:backend-presentation-format-matrix}.

\subsection{Evidence ladder for D3D11 3D draws}

The D3D11 3D checks form a progressively discriminating GPU-readback ladder. The first uses the
stride-16 \cnaclass{VertexPositionColor} path: a known blue clear was read back, an NDC-space
opaque-red triangle was then drawn through both indexed and non-indexed calls, and the same
region had to become red.  That proves the input layout, shaders, constant buffer, primitive
topology, rasterization, and readback path together; a stale clear or an unbound shader cannot
pass it accidentally.

The next rung separated texture sampling from vertex-color modulation.  A stride-20 textured
draw sampled a known texel exactly, through both indexed forms, while the stride-24 variant
multiplied a known vertex color through a white texture and read back that exact color.  The
stride-32 lighting path then exercises the per-light constant buffer: its unlit control is
byte-exact, while its lit control must differ from both the unlit result and the clear color.
That deliberately avoids pretending a hand-derived GPU floating-point result is more stable
than it is, while still distinguishing execution of the lighting branch.

Later tests applied the same discipline to features whose happy paths look deceptively similar.
Fog draws the same geometry with the flag off (exact vertex red) and at the fog end (exact fog
green); an environment-map fixture constrains the reflection to the cube's \texttt{-Z} face;
and a skinned fixture supplies an identity bone instead of an all-zero default matrix. These
fixtures falsify different failure modes. At the pin, the D3D11 smoke executable contains 147
checks exercised through Wine and DXVK GPU execution. The result supports the
specific paths named here; it is not a license to extrapolate that any arbitrary, untested
D3D11 state combination must therefore be correct.

\subsection{Viewport/scissor: three viewport paths, but only two scissor paths}

D3D9 and D3D11 take the simplest possible immediate route. Their public setters call the
corresponding native viewport and scissor setters; their rasterizer mapping separately controls
scissor enable. Neither SpriteBatch implementation replaces those native rectangles, and both
derive their 2D projection size from the currently-bound native viewport. D3D11 Smoke Check~O
reads the exact native viewport, depth range, and scissor back through the corresponding native
getters. That is strong native-state evidence but still one layer short of a public behavioral
pixel test because the check invokes renderer methods directly.

D3D9 implements both states, but its evidence label needed correction. The renderer plan
now calls viewport and scissor ``oracle-proven.'' A direct recount/search of the current
\XnaOracleSceneCount{}-scene \texttt{tools/xna-oracle/scenes} corpus finds cull-mode scenes but no viewport or
scissor scene, and the D3D9 smoke program's resize check proves a full-size post-reset viewport,
not a custom sub-region. The safest current claim is therefore native-source implementation
plus broad D3D9 draw coverage, not a zero-tolerance XNA comparison for these two states.

D3D12 now splits the two states. Its \texttt{SetViewport()} override stores rectangle and depth
range, and every ordinary, extended, instanced, and SpriteBatch command list obtains that value
through \texttt{GetEffectiveViewportEXT()}. The getter clamps depth to $[0,1]$ and falls back to
the current full target only when no positive custom rectangle is active. SpriteBatch also uses
the same effective width and height as its projection basis, so coordinates remain local to the
sub-viewport rather than being transformed twice. A native off-screen GPU check distinguishes
inside/outside pixels, two custom rectangles in successive draws, depth-range behavior, and the
full-target reset on a target switch. Scissor remains the negative half: there is still no
rectangle hook, \texttt{scissorTestEnable} is not part of the PSO, and every draw installs a
full-target \texttt{D3D12\_RECT}. The complete product comparison is
\S\ref{sec:backend-viewport-scissor-matrix}.

\subsection{Clear: exact native bits do not imply one shared public contract}

D3D11 immediately clears every bound RTV and selects the native depth and stencil clear bits
independently when a DSV exists. D3D12 mirrors that selection with synchronous command lists.
Both have strong direct-renderer proof: stencil-plane bytes are read back from the GPU, changed to a second
value, and depth is distinguished by whether one identical triangle passes. Those tests bypass
\cnaclass{GraphicsDevice}'s public route, but the route now derives depth and stencil availability
independently from the bound 2D/cube attachment instead of inferring stencil from every
non-\texttt{None} depth enum. D3D12 still requires a bound colour resource before its combination
helper can record any command list; once one is bound, however, a requested depth/stencil aspect
with no DSV is deliberately a legal no-op, matching D3D11.

D3D9's native \texttt{IDirect3DDevice9::Clear} bits are likewise immediate and individually
selected, but their final gate consults \texttt{depthStencilFormatOrdinal\_}, the presentation
format captured for the backbuffer. Binding a differently-formatted 2D or cube target never
updates that field. A target can therefore lose a depth clear when the backbuffer has no
depth, or receive a stencil bit selected because the backbuffer has one. The smoke test's seven
public combinations all use a D24S8 backbuffer and mostly check colour preservation, so they
cannot reveal this mismatch. See \S\ref{sec:backend-clear-matrix} for the corresponding
thirteen-product contrast within that audit cohort.

\section{MRT finalization across all bound targets}

An ordinary single \cnaclass{RenderTarget2D} owns a simple end-of-use rule. When it is unbound,
its \texttt{ResolveAnd\allowbreak{}Generate\allowbreak{}MipsEXT()} helper asks D3D11 to resolve the
multisampled draw attachment into the separate sampleable texture. If the target
has a mip chain, it then calls \texttt{GenerateMips()} on the shader-resource view. Thus the
texture a later sprite samples is the resolved, current image rather than the still-multisampled
attachment or an old mip level.

That rule was once silently absent for an \(N>1\) MRT bind. The single-target pointer could not
represent several targets, so the old \texttt{SetRenderTargets()} path bound all their RTVs with
\texttt{OMSetRenderTargets()} but did not call any target's finalization when the set was
replaced. A game could clear or draw to two MSAA targets successfully, then sample stale
single-sample resolve textures without an exception. The repair keeps a non-owning array plus a
count for the active MRT set. Both \texttt{SetRenderTargets()} and
\texttt{SetRenderTarget2D()} first run \texttt{Flush\allowbreak{}Pending\allowbreak{}MRT\allowbreak{}ResolveEXT()}, which visits every
stored target, invokes its \texttt{ResolveAnd\allowbreak{}Generate\allowbreak{}MipsEXT()}, clears the array entry, and only
then binds the next target or restores the back buffer. This matters for all three transitions:
MRT to back buffer, MRT directly to one target, and MRT to a different MRT set.

\texttt{D3D11\_Smoke} makes the error observable. It binds
two distinct \(8\times8\), 4x-MSAA targets, clears the shared MRT pass to
\((200,30,40,255)\), unbinds it, and reads both targets' \emph{sampleable} textures directly.
Both must contain that exact color. A repair that resolved only target zero, or that merely left
the MSAA draw attachment valid, fails the second readback. A second case binds another two-target
set, clears it to \((5,6,7,255)\), and switches straight to an unrelated single target rather
than unbinding to null; its first prior target must still read the new value. This distinguishes
all target transitions from a narrower null-unbind repair.

The shared helper also calls \texttt{GenerateMips()}, but the scope of its proof should remain
honest. The suite separately reads an \(8\times8\) single-target mipmapped render target's
\(4\times4\) and \(2\times2\) levels after a solid \((200,90,10,255)\) clear, proving that the
normal unbind path creates non-garbage downstream levels. The \emph{combination} of \(N>1\)
MRT and \texttt{mipMap=true} uses the same helper and is therefore wired, but is not independently
pixel-tested yet. That is a verification boundary, not a claim that the two mechanisms are
magically proven together.

\section{D3D12: shared foundations and split verification}

\texttt{D3D12} reuses D3D11's HLSL/DXBC bytecode, format/state mappings, and constant-buffer
layouts through \texttt{D3DCommon}. The implemented route includes a device, queues, descriptor
heaps, command lists, fences, per-resource barriers, pipeline-state objects, root signatures,
per-slot samplers, buffers, 2D/cube/3D textures, render targets with MRT, mip generation and
device-queried MSAA, occlusion queries, all ten stock shader variants, \cnaclass{SpriteBatch},
runtime-compiled sprite shaders, and device-removed recovery.

\begin{evidencenote}
The routine D3D12 CTest suite is entirely off-screen. A separate manual diagnostic, launched
through Proton, establishes swapchain creation and \texttt{Present()} through a window. That
manual result is not part of the routine automated suite because the Proton bootstrap is too
heavy for this development loop.
\end{evidencenote}

\subsection{Backbuffer readback is not part of the shared D3D11 foundation}

The repeated phrase ``GPU readback'' in this chapter names several different mechanisms.
D3D11 has a \texttt{ReadBackbuffer()} override: it always copies
\texttt{backBufferTexture\_} into a staging texture, honors \texttt{RowPitch}, and backs the
public \texttt{GraphicsDevice::GetBackBuffer\allowbreak Data()} route with exact windowed
clear/draw
tests. It does not redirect to a currently bound render target.

D3D12 has no corresponding override. A public backbuffer call reaches
\cnaclass{IGraphicsRenderer}'s default \texttt{runtime\_error} on a windowed device, while
\texttt{HeadlessEXT} has no swapchain to read in the first place. The D3D12 smoke suite's many
exact GPU pixel assertions call its working render-target-resource readback helpers directly;
the shared public SpriteBatch example below likewise reads the off-screen render target, not
\texttt{GetBackBufferData()}. Those results prove the named draw/resource path but cannot be
quoted as evidence for a public backbuffer bridge that does not exist. The complete renderer
matrix and FNA contrast are in \S\ref{sec:backend-readback-matrix}.

\subsection{Two frames in flight and blocking call sites}

Unlike the driver-managed D3D11 model, D3D12 makes the command allocator lifetime explicit.
CNA therefore creates exactly two \texttt{ID3D12CommandAllocator} objects
(\texttt{kFramesInFlight = 2}), one reusable graphics command list, one shared
\texttt{ID3D12Fence}, and one remembered fence value per allocator slot.  The essential
operation, \texttt{SignalAndWaitForFrameEXT(\allowbreak{}frameIndex)}, deliberately has a slightly
counter-intuitive ordering: it signals a \emph{new}, monotonically increasing value for the
frame just submitted, then waits only for the \emph{previous} value recorded for that same
slot before allowing its allocator to be reused.  Thus a call for slot~0 may submit value
\texttt{v2} while waiting for older \texttt{v0}; it must not wait for \texttt{v2} itself,
because doing so would erase the overlap the two slots exist to permit.

The smoke test invokes the primitive for slots 0, 1, and 0 again, verifies
that \texttt{v0 < v1 < v2}, and, on the second use of slot~0, requires
\texttt{GetCompletedValue()} to have reached \texttt{v0}.  It intentionally does \emph{not}
assert that \texttt{v2} has completed on return: \texttt{ID3D12CommandQueue::Signal()} is
asynchronous, so that would be a race, not a correctness condition.  A separate explicit
event wait establishes eventual completion. The test also places 48
64-MiB GPU copies ahead of the old fence and comparing that wait with an already-complete
control wait; the loaded wait is measurably slower, establishing that
\texttt{WaitForSingleObject()} supplies back-pressure.

The reusable two-slot primitive is directly tested, but the current ordinary clear/draw helper
uses \texttt{ExecuteCommandListAndWaitEXT()}, its intentionally simpler synchronous sibling:
it submits a closed list, signals a fresh fence value, and waits for \emph{that} value before
returning.  \texttt{Present()} likewise uses that synchronous helper for its explicit
back-buffer transition before calling \texttt{IDXGISwapChain3::Present()}.  Consequently a
``Two frames in flight'' therefore describes tested infrastructure and the allocator-reuse rule;
it does not claim that current ordinary CNA drawing overlaps two CPU frames. The off-screen test
paths retain a deterministic readback boundary.

Render-target,
depth-stencil, shader-resource, and sampler descriptors come from four separate heaps; their
current allocators are fixed-capacity, monotonically advancing bump allocators, not a general
free-list.  A destroyed resource does not return its descriptor slot during the device's
lifetime, and exhaustion throws a named \texttt{std::runtime\_error} instead of silently
aliasing another resource's view.  The capacities in the live renderer --- 64 RTV, 8 DSV,
64 CBV/SRV/UAV, and 16 sampler descriptors --- are therefore engineering limits of this
implementation, not Direct3D~12 limits. Destroying many transient resources can therefore
exhaust a descriptor heap even when GPU memory has been released; callers cannot assume
unlimited descriptor churn during one device lifetime.

\subsection{External Wine/vkd3d architecture mismatch}

Swap-chain creation under \emph{plain} Wine (as opposed to a Proton-managed launch) crashes
outright, reproduced twice with a full symbolized backtrace pointing to a null-pointer read
inside Wine's own \texttt{dxgi.dll}. The cause is external to
CNA: a mismatch between Debian's own system \texttt{dxgi.dll} and vkd3d-proton's separately
overridden \texttt{d3d12.dll} --- two DLLs that were never meant to be paired this way. Once
launched through Proton, which supplies the matched DLL pair expected by vkd3d-proton,
swapchain creation and a ten-frame clear-and-present loop through a window succeed in two runs.

\subsection{Explicitly corrected claims}

Earlier D3D12 documentation listed runtime state objects, sixteen-slot sampler state, public 2D
and cube render targets including MRT, \cnaclass{Texture3D}, and
\cnaclass{TextureCube::GetData()} as open. They are implemented at the pin. The live tracking
file remains the authority for revisions after this edition.

\section{One public \texttt{ShaderEffect} API, two Direct3D contracts}

Chapter~\ref{ch:shader-fx-gap} draws an important boundary: D3D11 and D3D12
compile custom HLSL at runtime, but their implementations are \cnaclass{SpriteBatch} custom
shader facilities, not general replacements for XNA's compiled \texttt{Effect}. The two
renderers accept the same constructor strings, compile the same \texttt{main} entry points as
\texttt{vs\_5\_0} and \texttt{ps\_5\_0}, and expose the same public setters. Under that shared
surface, both require one exact 32-byte vertex:

\begin{longtable}{@{}p{0.22\textwidth}p{0.18\textwidth}p{0.46\textwidth}@{}}
\toprule
\textbf{Semantic} & \textbf{Byte offset} & \textbf{Required HLSL type} \\
\midrule
\endhead
\texttt{POSITION0} & 0  & \texttt{float2}: pixel-space X and Y \\
\texttt{TEXCOORD0} & 8  & \texttt{float2}: texture U and V \\
\texttt{COLOR0}    & 16 & \texttt{float4}: sprite tint \\
\bottomrule
\end{longtable}

This is \cnaclass{SpriteBatch}'s own vertex, not an inferred convention. D3D11 creates a
three-element \texttt{ID3D11InputLayout} from the compiled vertex blob; D3D12 embeds the same
three elements in the custom effect's pipeline-state object. A shader declaring
\texttt{float3 POSITION}, a tangent, or a game-specific stride therefore cannot be made
compatible by supplying a matching \cnaclass{VertexDeclaration}: unlike EasyGL's separately
tested general-3D path, neither Direct3D renderer reads
\texttt{GpuDrawParams::customEffectRenderer} from its 3D draw dispatch.

\subsection{The names are documentation; the byte offsets are the API}

Both implementations store public uniform writes in the same 128-byte block. The caller's
\texttt{name} argument is deliberately ignored. What selects a value is which setter was
called:

\begin{longtable}{@{}p{0.19\textwidth}p{0.25\textwidth}p{0.48\textwidth}@{}}
\toprule
\textbf{Bytes} & \textbf{Writer} & \textbf{Meaning} \\
\midrule
\endhead
0--15 &
renderer only &
\texttt{vpSize}; \cnaclass{SpriteBatch} writes width and height before binding the effect. \\
16--79 &
\texttt{SetUniformMat4} &
one 64-byte matrix, regardless of the supplied name. \\
80--95 &
\texttt{SetUniformVec2/3/4} &
one overlapping vector slot; a later vector setter overwrites the components it supplies. \\
96--99 &
\texttt{SetUniformFloat/Int} &
one scalar slot. The integer overload stores a floating-point conversion, not integer bits. \\
100--127 &
none &
reserved padding in the current implementation. \\
\bottomrule
\end{longtable}

There are two further consequences that the common public surface does not make obvious.
\texttt{SetUniformFloatArray()} and \texttt{SetUniformVec2Array()} inherit
\cnaclass{IEffectRenderer}'s no-op defaults on both renderers. So do all three
\texttt{SetTexture()} renderer hooks. The sprite's primary resource still works because
\cnaclass{SpriteBatch} itself binds its texture and sampler to \texttt{t0}/\texttt{s0}; a
second 2D texture, cube map, or volume texture set through \cnaclass{ShaderEffect} does not.
Effect authors should therefore treat the layout above as a tiny fixed ABI, not as
name-reflected HLSL constants.

\subsection{Where D3D11 and D3D12 stop being the same}

D3D11 can bind the compiled vertex shader, pixel shader, input layout, and dynamic constant
buffer as separate context state. Every \texttt{Bind()} performs a write-discard map, copies
the 128 bytes, and attaches the buffer to \texttt{b0} for both shader stages. The native map
mode is \texttt{D3D11\_MAP\_WRITE\_DISCARD}.

D3D12 has no equivalent independent shader bind. Its effect constructor must build a complete
\texttt{ID3D12PipelineState} up front, using the same cached root-signature shape as the stock
sprite path: one CBV at \texttt{b0}, one SRV table at \texttt{t0}, and one sampler table at
\texttt{s0}. The 128-byte constant buffer lives in a persistently mapped upload resource;
\texttt{Bind()} only copies the current values, while
\texttt{D3D12SpriteBatchRenderer::FlushBatch()} binds the resource and the prebuilt PSO inside
its command list.

That PSO also makes D3D12's boundary stricter and directly visible in source: it is compiled
for one \texttt{DXGI\_FORMAT\_R8G8B8A8\_UNORM}, single-sample color target, with blending,
depth, and stencil disabled and culling set to none. Those choices match the custom-effect
pixel test. They do not form a dynamic promise that the same object can be reused for an MRT
set, an MSAA target, another color format, or caller-selected pipeline state; those dimensions
would require separate PSOs and a cache key that the current effect renderer does not build.

\subsection{Shared D3D11/D3D12 pixel-test example}

The two smoke tests use byte-for-byte equivalent HLSL to turn a solid red texture into exact
cyan. The vertex shader's \texttt{vpSize} field is at byte zero because the renderer fills it;
the four padding vectors advance the next public vector slot to byte 80:

\begin{lstlisting}
const char* vertexHlsl = R"(
struct VSIn  { float2 pos : POSITION0;
               float2 uv  : TEXCOORD0;
               float4 col : COLOR0; };
struct VSOut { float4 pos : SV_Position;
               float2 uv  : TEXCOORD0;
               float4 col : TEXCOORD1; };
cbuffer CB : register(b0) {
    float4 vpSize;
    float4 pad1[4];
    float4 uColor;
    float4 uFloat0;
};
VSOut main(VSIn input) {
    VSOut output;
    float2 ndc = (input.pos / vpSize.xy) * 2.0 - 1.0;
    output.pos = float4(ndc.x, -ndc.y, 0.0, 1.0);
    output.uv  = input.uv;
    output.col = input.col;
    return output;
})";

const char* pixelHlsl = R"(
Texture2D texSampler : register(t0);
SamplerState texSamplerSampler : register(s0);
struct PSIn { float4 pos : SV_Position;
              float2 uv  : TEXCOORD0;
              float4 col : TEXCOORD1; };
float4 main(PSIn input) : SV_Target {
    float4 source = texSampler.Sample(texSamplerSampler, input.uv);
    return float4(1.0 - source.rgb, 1.0);
})";

ShaderEffect invert(device, vertexHlsl, pixelHlsl);
if (!invert.IsEffectValid()) {
    throw std::runtime_error("custom HLSL did not compile");
}

SamplerState pointClamp = SamplerState::PointClamp;
SpriteBatch batch(device);
batch.Begin(SpriteSortMode::Deferred, BlendState::Opaque,
            &pointClamp, nullptr, nullptr, &invert);
batch.Draw(redTexture, destination, Color::White);
batch.End();
\end{lstlisting}

D3D11 runs this through a windowed Wine/DXVK device and reads the sprite's top-left region.
D3D12 runs the same public \cnaclass{GraphicsDevice}/\cnaclass{SpriteBatch} sequence against a
\texttt{HeadlessEXT} off-screen device, then reads the render target. In both cases solid
\((255,0,0,255)\) must become exactly \((0,255,255,255)\). That result proves more than
successful compilation: the fixed input layout, automatic viewport slot, \texttt{b0}
constant-buffer/root binding, \texttt{t0}/\texttt{s0} resource path, custom pixel shader, and
GPU readback all participated in one observable public-API draw.

\section{Persistent Direct3D resources implement the shared usage contract}

All three Direct3D factories receive the preserve Boolean for 2D and cube targets but need not
translate it into a native load action: their resources persist across binding. Shared code
supplies the visible policy. \texttt{DiscardContents} clears black plus each available depth
and stencil aspect; \texttt{PreserveContents} and \texttt{PlatformContents} do not. The first
descriptor controls the same policy for MRT. A redundant bind is not a no-op: native
resolve/finalization can run before the resource is rebound, and a discard target receives its
three-aspect clear again.

Public pixel fixtures establish this policy. The shared
cube-usage fixture renders an asymmetric face, rebinds it and adds only a small marker, so a
discard masquerading as preservation cannot pass. Parallel depth and stencil fixtures preserve
an occluder or stamp across a bind cycle, while pass-boundary and ordered-clear tests separate
usage policy from explicit clear chronology. D3D9 and D3D11 register these tests in their routine
renderer suites; D3D12 cross-builds the same public fixtures, while its headless/direct smoke
coverage supplies the local GPU mechanisms. See \S\ref{sec:backend-rendertarget-usage-matrix}.

\section{Cube-target MSAA, mip generation, and public finalization}

Shared \texttt{D3DCommon} tables do not imply identical resource algorithms. D3D11 and D3D12
use parallel designs for \cnaclass{RenderTargetCube} MSAA but different mip-generation routes.
An interface-default audit also exposed a public routing hole after those mechanisms landed:
early smoke checks invoked \texttt{BindAsRenderTargetFace()} and
\texttt{UnbindAsRenderTarget()} directly, while the public route never finalized the outgoing
cube. That defect is now closed. The following subsections describe both the renderer algorithms
and the tracking edge that connects them to the public device route.

\subsection{\cnaclass{RenderTargetCube} MSAA}

Neither \texttt{D3D11\_RESOURCE\_MISC\_TEXTURECUBE} nor
a \texttt{D3D11\_SRV\_DIMENSION\_TEXTURECUBE} view can ever be multisampled, and \texttt{D3D12}'s
equivalent \texttt{D3D12\_SRV\_DIMENSION\_TEXTURECUBE} has no multisampled variant either --- so
on both renderers, the MSAA color resource becomes a plain, non-cube six-slice
\texttt{Texture2DMSArray} used only as a render-target view, while a second, single-sample,
cube-flagged resource receives \texttt{ResolveSubresource()} on unbind. The shader-resource view
targets the resolved resource. D3D12 additionally requires two explicit
resource-barrier transitions (\texttt{RESOLVE\_SOURCE} on the color resource,
\texttt{RESOLVE\_DEST} on the resolve target) that \texttt{D3D11} does not need at all --- its
context-level \texttt{ResolveSubresource()} call handles the equivalent state transition
implicitly. Both renderers resolve only the currently active face on unbind, the same
one-face-at-a-time convention their mip-chain handling follows. On both renderers, an
$8\times8$ cube face requested at 4x MSAA is cleared, resolved, and read back exactly;
\texttt{GetMultiSampleCount()} must also report four.

\subsection{Mip-chain generation: one native GPU call versus an explicit CPU round trip}

After drawing one cube face, D3D11 calls
\texttt{ID3D11DeviceContext::GenerateMips()} on a shader-resource view spanning all six faces.
The GPU regenerates every face's chain; unchanged level-zero faces reproduce their prior mips.
D3D12 has no corresponding convenience operation. Its
\texttt{GenerateMipsEXT()} reads the active face's preceding level through
\texttt{ReadbackSubresourceRGBA8()}, box-filters it on the CPU, and uploads the next level with
\texttt{UploadSubresourceRGBA8()}. It repeats this round trip for each level of that face.
Both routes produce the required chain, but their synchronization and cost differ substantially.
Shared bytecode and state mappings do not imply a shared resource algorithm.

\subsection{How renderer tracking closes the public finalization edge}

D3D11 and D3D12 now override the cube-face route and retain the active cube separately from the
single-2D and MRT state. Before any cube-to-cube, cube-to-2D, cube-to-MRT or cube-to-backbuffer
transition, \texttt{FlushPendingCubeResolveEXT()} clears that tracking pointer and calls the
outgoing cube's \texttt{UnbindAsRenderTarget()}. That ordering is important: the callback performs
the active-face MSAA resolve and mip generation before the new destination replaces native state.
The plural normalized-descriptor path recognizes a single cube face and delegates to this same
route; cube faces in a multi-target set remain an explicit renderer rejection rather than being
silently treated as 2D targets.

The original direct smoke checks still prove the two low-level algorithms described above, but
they are no longer the only evidence. Shared \cnaclass{GraphicsDevice} fixtures now drive an
asymmetric rendered cube face through public binding and readback, rebind it for the usage test,
and isolate preserved multisample content across two faces. Those fixtures make a skipped
finalization, wrong face/subresource, stale resolve or false preservation visibly fail. Thus the
former caution remains a useful account of how the hole escaped an implementation-local smoke
test, not a current feature limitation.

\section{\texttt{HeadlessEXT}: routine off-screen D3D12 coverage}

\cnaclass{SpriteFont}, \cnaclass{Model}, and part of the
\texttt{Texture2D.FromStream}/\texttt{SaveAsPng} route require a \cnaclass{GraphicsDevice}.
That constructor formerly created an SDL window for every renderer except \texttt{HEADLESS} and
\texttt{SOFTWARE}. On D3D12, window creation also required the Proton-managed swapchain route,
which was too expensive for the routine plain-Wine test loop.

\cnaclass{PresentationParameters} therefore gained the CNA-only \texttt{HeadlessEXT} property:

\begin{lstlisting}
CNAEXT bool getHeadlessEXTProperty() const;
CNAEXT void setHeadlessEXTProperty(bool value);
\end{lstlisting}

The property defaults to \texttt{false}; only D3D12 honors it at the pin. D3D11 always creates a
swapchain and EasyGL requires a window-bound GL context, so both reject it. A headless D3D12
device renders only to explicit targets and cannot present. This supports routine tests as well
as off-screen applications such as server-side rendering and thumbnail generation.

\subsection{Four D3D12 defects exposed by headless GraphicsDevice tests}

The first routine public-device tests exposed four corrected D3D12 defects:

\begin{itemize}
  \item \cnaclass{SpriteBatch} ignored sampler filter and address-mode updates, inheriting state
        from a prior 3D draw;
  \item render targets failed to bind their depth-stencil views;
  \item the six combined color/depth/stencil clear variants threw; and
  \item \texttt{SetDepthTestEnabled}, \texttt{SetDepthWriteEnabled}, and
        \texttt{SetBlendEnabled} threw when the \cnaclass{Model} route called them.
\end{itemize}

All four are fixed at the pin. Their discovery also shows why direct renderer smoke tests and
public \cnaclass{GraphicsDevice} fixtures are complementary: the former had not traversed these
shared-framework call paths.
